% Self-contained proof module. Citation keys: James1978, Quillen1973.
\section{Regular local coefficient rings}

The positive result over a discrete valuation ring extends to every
commutative noetherian regular local ring of characteristic zero.  The
argument uses ordinary and modular representation theory only to identify
two triangular systems of generators; the passage between the fibres is
an argument with finite resolutions.

\begin{theorem}\label{thm:regular-local}
Let $A$ be a commutative noetherian regular local ring of characteristic
zero.  For every $n\geq 0$, extension of scalars to a characteristic-zero
field containing $A$ induces an isomorphism
\[
 K_0(\operatorname{Lat}_A(S_n))\simeq R_{\mathbb Q}(S_n).
\]
Consequently the homomorphism
\[
 \mathbb Z[\rho_1,\rho_2,\ldots]\longrightarrow
 \bigoplus_{n\geq 0}K_0(\operatorname{Lat}_A(S_n)),
\]
with multiplication given by induction and $\rho_n$ the trivial
$A[S_n]$-lattice, is an isomorphism.
\end{theorem}

\begin{proof}
Fix $n$ and write $G=S_n$.  We first compare lattices with arbitrary
finitely generated modules.  Put $d=\dim A$.  Any finitely generated
$A[G]$-module $N$ has a resolution
\[
 0\longrightarrow L\longrightarrow A[G]^{a_{d-1}}
 \longrightarrow\cdots\longrightarrow A[G]^{a_0}
 \longrightarrow N\longrightarrow 0
\]
in which $L$ is finite free over $A$ (if $d=0$, take $L=N$): choose successive finite free
$A[G]$-covers, and use $\operatorname{gl.dim}(A)=d$ for the last kernel.
The category of lattices is closed under extensions and under kernels
of surjections between its objects, since such surjections split over
$A$.  The resolution theorem therefore gives
\begin{equation}\label{eq:lattice-G0}
 K_0(\operatorname{Lat}_A(G))\simeq G_0(A[G]).
\end{equation}
Here $G_0$ denotes the Grothendieck group of finitely generated modules.
At the level used here, the inverse sends a module to the alternating
sum of any finite lattice resolution.  Comparison of resolutions and
the horseshoe construction prove independence and additivity; an
acyclic bounded complex of $A$-free modules has $A$-free cycles, by
successively splitting from its right end.  This is also the
$K_0$ case of the resolution theorem in
\cite[\S4, Theorem~3 and Corollary~1]{Quillen1973}.

If the residue field of $A$ has characteristic zero, every positive
integer is invertible in $A$.  The split Wedderburn decomposition
\[
 \mathbb Q[S_n]\simeq
 \prod_{\lambda\vdash n}\operatorname{Mat}_{d_\lambda}(\mathbb Q)
\]
then gives $A[S_n]\simeq\prod_{\lambda\vdash n}
\operatorname{Mat}_{d_\lambda}(A)$.  Morita equivalence and
$G_0(A)=\mathbb Z[A]$ prove the assertion in this case.

Suppose instead that the residue characteristic is $p>0$, and put
$B=A/pA$ and $C=A[1/p]$.  Every prime different from $p$ is already a
unit in $A$, so $C$ is a $\mathbb Q$-algebra.  Localization of
finitely generated modules gives an exact sequence
\begin{equation}\label{eq:regular-localization}
 G_0(\mathcal T)\longrightarrow G_0(A[G])
 \longrightarrow G_0(C[G])\longrightarrow 0,
\end{equation}
where $\mathcal T$ consists of modules killed by a power of $p$.
Filtering a module by powers of $p$ shows that the image of the first
map is generated by the classes of finitely generated $B[G]$-modules.
We prove that every such class is zero in $G_0(A[G])$.

Let $J$ be the radical of the finite-dimensional algebra
$\mathbb F_p[G]$.  By the splitting-field theorem for symmetric groups,
\[
 \mathbb F_p[G]/J\simeq
 \prod_{\mu\ p\text{-regular}}\operatorname{Mat}_{e_\mu}(\mathbb F_p).
\]
The ideal $B\otimes_{\mathbb F_p}J$ is nilpotent in $B[G]$.
Its powers filter every finitely generated $B[G]$-module.  Morita
equivalence for the displayed product shows that each successive
quotient is a direct sum of modules
\[
 M\otimes_{\mathbb F_p}D^\mu,
\]
where $M$ is a finitely generated $B$-module and $D^\mu$ is a simple
$\mathbb F_p[G]$-module.  The group acts trivially on $M$.

For a partition $\nu$ of $n$, let $Y^\nu_{\mathbb Z}$ denote the
permutation lattice on the cosets of its Young subgroup.  James's
splitting and triangularity theorems
\cite[Theorem~11.5, p.~40, and Theorem~12.1, p.~42]{James1978}
imply that the classes of $Y^\nu_{\mathbb F_p}$, with $\nu$
$p$-regular, form a $\mathbb Z$-basis of
$G_0(\mathbb F_p[G])$.  Indeed, their composition multiplicities
against the $D^\mu$, ordered by dominance, form a unitriangular
matrix.  Write
\[
 [D^\mu]=\sum_\nu a_{\mu\nu}[Y^\nu_{\mathbb F_p}].
\]
Tensoring an exact sequence of $\mathbb F_p[G]$-modules with the
$\mathbb F_p$-vector space $M$ is exact.  Therefore
\begin{equation}\label{eq:kill-modular-class}
 [M\otimes_{\mathbb F_p}D^\mu]
 =\sum_\nu a_{\mu\nu}
 [M\otimes_A Y^\nu_A]
 \quad\hbox{in }G_0(A[G]).
\end{equation}
The finitely generated $A$-module $M$ has a finite free resolution
$F_\bullet$.  Tensoring this resolution with the $A$-free lattice
$Y^\nu_A$ gives
\[
 [M\otimes_A Y^\nu_A]
 =\left(\sum_i(-1)^i\operatorname{rank}_A F_i\right)
 [Y^\nu_A]=0.
\]
The last equality holds because $pM=0$ and $p\ne0$ in the domain
$A$, so $M$ has rank zero.  Equation~\eqref{eq:kill-modular-class}
proves the required vanishing.  Thus
\begin{equation}\label{eq:regular-invert-p}
 G_0(A[G])\simeq G_0(C[G]).
\end{equation}

Finally, localization gives a surjection $G_0(A)\to G_0(C)$.
Since $G_0(A)=\mathbb Z[A]$, and rank over the fraction field sends
$[C]$ to $1$, this map is an isomorphism and
$G_0(C)=\mathbb Z[C]$.  Tensoring the rational split Wedderburn
decomposition with $C$ now yields
\[
 G_0(C[G])\simeq\bigoplus_{\lambda\vdash n}G_0(C)
 \simeq\mathbb Z^{\{\lambda\vdash n\}}.
\]
These identifications agree with extension to the fraction field of
$A$ and hence with the ordinary representation group of $S_n$.
The ordinary Young permutation classes form a unitriangular integral
basis, and the monomial $\rho_{\lambda_1}\cdots\rho_{\lambda_r}$
is precisely the class induced from the corresponding Young subgroup.
The degreewise isomorphisms are compatible with induction, proving
the final assertion.
\end{proof}

\begin{remark}
No regularity of $A/pA$ was used.  In particular, the theorem includes
ramified regular local rings.  Regularity of $A$ enters in the
comparison with $G_0$ and in the finite free resolution of the
multiplicity module $M$.
\end{remark}
